\chapter{Потоки й обертання плазми та радіальне електричне поле}\label{ch:flows}

Це центральний розділ посібника. Ланцюжок такий: рівняння руху іонів $\to$ радіальний баланс сил і $\Er$ (п.~\ref{sec:06-force}) $\to$ структура потоку на магнітній поверхні (п.~\ref{sec:06-surface}) $\to$ неокласичний полоїдальний потік (п.~\ref{sec:06-pol}) $\to$ транспорт тороїдального імпульсу (пп.~\ref{sec:06-tor}--\ref{sec:06-intrinsic}) $\to$ зворотний зв'язок через шир $\ExB$ (п.~\ref{sec:06-shear}). Далі розглянуто край (п.~\ref{sec:06-sol}), діагностики (п.~\ref{sec:06-meas}) і замкнену модель для проєкту (п.~\ref{sec:06-model}). Зіткнення, режими ν* і 1.5D-рівняння для $n$, $T$, $\psi$ викладено в розд.~\ref{ch:transport}, обертову систему відліку і потенціал $\Phi_1$~--- в розд.~\ref{ch:particles}. Швидкість сорту $s$ позначаємо $\uvec_s$ (у джерелах $U$, $V$, $v$), полоїдальний кут~--- $\theta$ (у~\cite{Singh2004,Seto2014} його позначено грецькою «хі»), тороїдальний~--- $\varphi$ (у~\cite{Seto2014}~--- «дзета»).

\section{Радіальний баланс сил і поле $\Er$}\label{sec:06-force}

Рівняння руху сорту $a$ у двовимірній моделі Сето--Фукуями~\cite{Seto2014} (СІ):
\begin{equation}\label{eq:06-mom}
  \pa_t(m_an_a\uvec_a)+\grad\cdot\tens{P}_a=e_an_a(\vect{E}+\uvec_a\times\Bvec)+\vect{F}_a+\vect{F}^{\mathrm{QL}}_a+\vect{S}_{ma},
\end{equation}
де $\tens{P}_a=p_a\tens{I}+m_an_a\uvec_a\uvec_a+\visc_a$~--- повний тензор напружень, $\vect{F}_a$~--- тертя, $\vect{F}^{\mathrm{QL}}_a$~--- квазілінійна сила турбулентності, $\vect{S}_{ma}$~--- джерело імпульсу. Упорядкуємо за $\delta=\rho_a/L_\perp\ll1$. Відношення інерції $m_an_a(\uvec\cdot\grad)\uvec$ до сили Лоренца $e_an_a\uvec\times\Bvec$ має порядок $u/(\Omc{a}L)\lesssim\delta$ навіть при $u\sim v_{\mathrm{th}}$, в'язкість і тертя ще менші. Тому в перпендикулярному напрямку лишається рівновага нульового порядку. Її радіальну проєкцію Сето і Фукуяма записують у довільних координатах магнітних поверхонь~\cite[рівн.~46]{Seto2014}:
\begin{equation}\label{eq:06-seto}
  \grad\rho\cdot\grad p_a=e_an_a\bigl[E_\rho+\grad\rho\cdot(\uvec_a\times\Bvec)\bigr].
\end{equation}
Нехай $\grad\rho\parallel\hat{\vect{r}}$, а трійка $(\hat{\vect{r}},\hat{\vect{\theta}},\hat{\vect{\varphi}})$ права. Тоді $(\uvec\times\Bvec)_r=u_\theta\Btor-\utor\Bpol$, і для основних іонів ($e_a=Z_ie$)
\begin{equation}\label{eq:06-Er}
  \boxed{\;\Er=\frac{1}{Z_ien_i}\frac{\pa p_i}{\pa r}-u_{\theta i}\Btor+u_{\varphi i}\Bpol\;}
\end{equation}
Саме це співвідношення використовують Singh та ін.~\cite[рівн.~2]{Singh2004} (при $Z_i=1$ перший доданок дорівнює $(T_i/e)\pa_r\ln p_i$) та Ida~\cite{Ida2001}. Поле $\Er$ одне для всіх сортів, але його «склад» у кожного сорту свій: у важкої домішки ($Z_I\gg1$) доданок тиску малий, і $\Er$ майже повністю визначають її швидкості. Саме тому $\Er$ відновлюють зі спектроскопії домішок (п.~\ref{sec:06-meas}).

\textbf{Знаки.} \eqref{eq:06-Er} правильна лише у правій трійці $(r,\theta,\varphi)$. Якщо $\varphi$ узгоджено з циліндричними $(R,\varphi,Z)$, то $\hat{\vect{\theta}}=\hat{\vect{\varphi}}\times\hat{\vect{r}}$ на зовнішньому обводі дорівнює $-\hat{\vect{Z}}$, тобто θ зростає \emph{вниз}. Якщо відлічувати θ вгору, обидва доданки зі швидкостями змінюють знак. Форма~\eqref{eq:06-seto} від вибору орієнтації не залежить. На практиці напрямки задають фізично: «co»~--- паралельно $\Ip$, полоїдальний~--- «в електронному (іонному) діамагнітному напрямку». Дрейф $\ExB$ при $\Er<0$ іде в електронному діамагнітному напрямку, тому полоїдальний потік у цьому напрямку відповідає $\Er<0$~\cite{Ida2001}.

\textbf{Три граничні випадки.} (1)~Ядро з NBI (TFTR, практикум): домінує $\utor\Bpol$. (2)~Геліотрон CHS: доданки тиску і $\utor\Bpol$ малі, і $\Er\approx u_\theta\Btor$; поле з балансу сил узгоджується з прямим виміром HIBP~\cite{Ida2001}. (3)~Край токамака: усі три доданки однакового порядку (задача~2). Отже, \eqref{eq:06-Er} сама нічого не передбачає, поки не відомі $u_\theta$ і $\utor$. Їх дають пп.~\ref{sec:06-pol}--\ref{sec:06-intrinsic}. У двовимірній моделі $E_\rho$ визначають із закону Гаусса, а не з квазінейтральності~\cite{Seto2014}.

\section{Потік на магнітній поверхні}\label{sec:06-surface}

\begin{derivation}[форма потоку нульового порядку; EFIT, $\Bvec=\grad\psi\times\grad\varphi+F\grad\varphi$]
У нульовому порядку $p_i=p_i(\psi)$ і $\Phi=\Phi(\psi)$ (розд.~\ref{ch:equilibrium}). З \eqref{eq:06-mom} маємо $\uvec_\perp=\vE+\frac{\Bvec\times\grad p_i}{Z_ien_iB^2}$, тобто обидва дрейфи пропорційні $\Bvec\times\grad\psi/B^2$. Розпишемо векторний добуток: $\Bvec\times\grad\psi=F\grad\varphi\times\grad\psi+(\grad\psi\times\grad\varphi)\times\grad\psi=-F(\Bvec-F\grad\varphi)+|\grad\psi|^2\grad\varphi$. Оскільки $R^2B^2=F^2+|\grad\psi|^2$,
\begin{gather*}
  \frac{\Bvec\times\grad\psi}{B^2}=R^2\grad\varphi-\frac{F}{B}\unitb
  \quad\Rightarrow\quad
  \uvec_{\perp}=\omega_\psi(\psi)\Bigl(R^2\grad\varphi-\frac{F}{B}\unitb\Bigr),\\
  \omega_\psi\equiv\omega_E+\frac{1}{Z_ien_i}\frac{\dd p_i}{\dd\psi},\qquad \omega_E\equiv\frac{\dd\Phi}{\dd\psi}.
\end{gather*}
Додамо $\upar\unitb$. Доданок $\omega_\psi R^2\grad\varphi$ осесиметричний і бездивергентний, тому стаціонарна неперервність $\grad\cdot(n_i\uvec_i)=0$ зводиться до $\Bvec\cdot\grad\bigl[n_i(\upar-\omega_\psi F/B)/B\bigr]=0$. Отже,
\begin{equation}\label{eq:06-uform}
  n_i\uvec_i=n_i\,\omega_\psi(\psi)\,R^2\grad\varphi+K_i(\psi)\,\Bvec,\qquad
  u_{\theta i}=\frac{K_i\Bpol}{n_i},\qquad u_{\varphi i}=\omega_\psi R+\frac{K_i\Btor}{n_i}.
\end{equation}
\end{derivation}

Отже, потік на поверхні задають \emph{дві} функції потоку: частота твердотільного обертання $\omega_\psi$ і полоїдальна функція $K_i$. Полоїдальна швидкість змінюється по поверхні як $\Bpol/n_i$, тому вздовж обводу вона не стала. Радіальне поле $|\Er|=|\omega_E|\,|\grad\psi|=|\omega_E|R\Bpol$ (ψ на радіан; для повного потоку $2\pi\psi$ частота в $2\pi$ разів менша). У~\cite[рівн.~53]{Seto2014} ту саму форму записано як $\bar{\uvec}_a=\omega_{ua}R^2\grad\varphi+L_{ua}\Bvec$, тобто $L_{ua}=K_a/n_a$ (у них $\Bvec=\grad\varphi\times\grad\psi+F\grad\varphi$, знак ψ протилежний).

\textbf{Модель Сігмара--Заніно--Хсу.} Нехай NBI розганяє плазму до $V_\varphi\sim v_{\mathrm{th}}$. Тоді домінує $\vE$, а діамагнітний дрейф вищого порядку. Автори записують~\cite[с.~7--8, рівн.~8--10]{Sigmar1987} (переведено з СГС)
\begin{equation}\label{eq:06-sigmar}
  \uvec^{(0)}=\Bigl(\upar+\omega_{\mathrm S}\frac{F}{B}\Bigr)\unitb-\omega_{\mathrm S}R^2\grad\varphi,\qquad
  \omega_{\mathrm S}=\frac{\pa\Phi}{\pa\psi_{\mathrm S}}\quad\Bigl(\text{СГС: }c\,\frac{\pa\Phi}{\pa\psi}\Bigr),
\end{equation}
з $\Bvec=\grad\varphi\times\grad\psi_{\mathrm S}+F\grad\varphi$, тобто $\psi_{\mathrm S}=-\psi$ і $\omega_{\mathrm S}=-\omega_E$. Отже, \eqref{eq:06-sigmar} збігається з \eqref{eq:06-uform} при $\omega_\psi\to\omega_E$ (додаток~\ref{app:cgs}). Неперервність на поверхні набуває вигляду $\pa_tn_j+\Bvec\cdot\grad\bigl[n_j(u_{\parallel j}-\omega_EF/B)/B\bigr]=0$. Отже, полоїдальний потік $V_p=(\upar+\omega_{\mathrm S}F/B)\Bpol/B$ перерозподіляє густину $n_j(\psi,\theta)$, і в стаціонарі $n_jV_{pj}/\Bpol$ стає функцією потоку (це і є $K_j$)~\cite{Sigmar1987}. Для домішок із сильною полоїдальною асиметрією (ASDEX, амплітуди Фур'є $\order{1}$) це основний механізм.

Радіальна проєкція рівняння руху, підсумована по сортах, з законом Ампера $\pa_t(\omega_{\mathrm S}|\grad\psi|^2)=\vect{J}\cdot\grad\psi_{\mathrm S}/\varepsilon_0$ після усереднення дає рівняння для $\omega\propto\Er$~\cite[с.~9--11, рівн.~13--14]{Sigmar1987}:
\begin{equation}\label{eq:06-omega}
  \frac{\pa\omega_{\mathrm S}}{\pa t}+C_1\omega_{\mathrm S}^2+C_2\omega_{\mathrm S}+C_3=0,\qquad
  D=\sum_j\fsa{m_jn_j}+\varepsilon_0\fsa{B^2},\quad C_1=\frac{1}{D}\frac{\Btor}{B}\Bigl\langle\sum_j m_jn_j\,\unitb\cdot\grad R\Bigr\rangle .
\end{equation}
$C_1$ походить від відцентрової інерції, $C_2$~--- від гальмування $\fsa{mn\nu_d}$, від $\upar$ і від паралельної в'язкості $\eta_0$ (з множником $\Btor^2/\Bpol^2$), $C_3$~--- від джерел: тороїдальної проєкції $M_{\perp\varphi}$ та, непрямо, від $M_\parallel$ через $\upar^2$. Тому навіть суто паралельна інжекція створює велике $\Er$, коли $\upar\sim v_{\mathrm{th}\,i}$. Доданок $\varepsilon_0\fsa{B^2}$ (у СГС $\fsa{B^2}/4\pi c^2$) відносно малий, як $v_{\mathrm A}^2/c^2$. Квадратичність \eqref{eq:06-omega} дає два стаціонарні корені $\omega_\pm=[-C_2\pm\sqrt{\Delta}]/(2C_1)$, $\Delta=C_2^2-4C_1C_3$, тобто можливу біфуркацію~\cite{Sigmar1987}. Лінеаризація (наш аналіз) дає $\delta\dot\omega=-(2C_1\omega_\pm+C_2)\,\delta\omega=\mp\sqrt\Delta\,\delta\omega$: стійкий лише корінь $\omega_+$, а час релаксації дорівнює $1/\sqrt\Delta$. При $\Delta\to0$ корені зливаються (складка, пор. неєдиність розв'язків GS у розд.~\ref{ch:equilibrium}).

\textbf{Масштаби часу.} Паралельна в'язкість гасить полоїдальний рух за час $\tau_p$, а на повільному масштабі $\tau_T$ лишається тороїдальне обертання. Це загальна властивість осесиметрії: усереднена по поверхні тороїдальна проєкція паралельної в'язкої сили дорівнює нулю~\cite[рівн.~42]{Seto2014}, тоді як полоїдальна ні. У простій моделі Хіршмана $\Er=R\Bpol\omega$ (переведено з СГС) релаксує з тим самим темпом $\tau_p^{-1}$, що й полоїдальний потік~\cite[дод.~C]{Sigmar1987}. Сігмар зі співавторами наголошують, що в експериментах $\tau_p$ і $\tau_T$ розділені значно гірше, ніж дає класична поперечна в'язкість $\eta_1\sim\eta_0(\Omc{i}\tau_i)^{-2}$: $\tau_T^{-1}(\text{експ.})\sim(50\text{--}100)\,\eta_1/(m_in_ia^2)$~\cite[с.~11]{Sigmar1987} (у джерелі скорочено $\eta_1/a^2$). Тому гальмування моделюють феноменологічно, $-m_jn_j\nu_{d}\uvec_j$. Аномальний перенос імпульсу розглянуто в п.~\ref{sec:06-intrinsic}.

\section{Неокласичний полоїдальний потік}\label{sec:06-pol}

Паралельна в'язкість «закріплює» $K_i$: полоїдальна швидкість перестає бути вільним параметром і виражається через $\pa_rT_i$. Singh та ін. наводять стандартний неокласичний результат~\cite[рівн.~2]{Singh2004}
\begin{equation}\label{eq:06-utheta}
  u_{\theta i}=\mathcal{K}\,\frac{1}{e\Btor}\frac{\pa T_i}{\pa r},\qquad
  \mathcal{K}\approx1{,}17\ (\text{банановий}),\quad -0{,}50\ (\text{плато}),\quad -1{,}83\ (\text{Пфірша--Шлютера}).
\end{equation}
Підставивши \eqref{eq:06-utheta} у \eqref{eq:06-Er} ($Z_i=1$), отримаємо
\begin{equation}\label{eq:06-Erneo}
  \Er=u_{\varphi i}\Bpol+\frac{T_i}{e}\Bigl[\frac{\pa\ln n_i}{\pa r}+(1-\mathcal{K})\frac{\pa\ln T_i}{\pa r}\Bigr].
\end{equation}
Цю форму використовують у~\cite{Singh2004}. Локальний шир $\ExB$ можна змінювати або градієнтом тиску, або тороїдальним обертанням.

\begin{corpusexample}[як отримати $-1{,}83$ у рідинній моделі]
У режимі Пфірша--Шлютера Singh та ін. беруть дворідинні рівняння Брагінського з поправками Михайловського--Ципіна до тензора в'язкості~\cite[дод.~A--B]{Singh2004}. Усереднення паралельної компоненти рівняння руху з вагою $\mathcal{J}B$ дає $\fsa{\pa_\theta\ln B\;\pi_{0\parallel}}=0$. Туди входять полоїдальні гармоніки $b^{(1)}$, $n^{(1)}$, $t^{(1)}$ і в'язкість $\eta_{0i}=0{,}96\,p_i/\nu_i$. Без поправки Михайловського--Ципіна з цієї умови випливає $u_{\theta i}=0$. З поправкою отримуємо $u_{\theta i}=-1{,}84\,\frac{T_i}{e\Bzero}\pa_r\ln T_i$ (рівн.~B10), що близько до кінетичного результату Хінтона--Хейзелтайна. В основному тексті статті стоїть $-1{,}83$. Розбіжність у третьому знаку є внутрішньою неузгодженістю джерела.
\end{corpusexample}

\begin{outofcorpus}[банановий режим]
При $\nustar\ll1$ паралельну в'язкість визначають захоплені частинки: $\fsa{\Bvec\cdot\grad\cdot\visc_i}\approx3\fsa{(\unitb\cdot\grad B)^2}\mu_i\,u_{\theta i}/\Bpol$, $\mu_i\propto f_t\nu_i$ (див. п.~\ref{sec:05-closure}). Баланс з тертям і тепловою силою дає $\mathcal{K}\approx1{,}17$ для чистої плазми, $\varepsilon\ll1$. Для JET у ядрі потрібен саме цей режим, а на краю~--- плато і Пфірша--Шлютера. Коефіцієнт $1{,}17$ є в корпусі лише як число~\cite{Singh2004}; виведення $1{,}17$ поза корпусом (Хіршман--Сігмар 1981).
\end{outofcorpus}

\textbf{Багатосортове узагальнення.} Нехай $h_{a\parallel}=2Q_{a\parallel}/5p_a$ (як $\vect{h}_s$ у розд.~\ref{ch:transport}), $\fsa{u_{a\parallel}B}=V_{1a}B+L_{ua}\fsa{B^2}$ і $\fsa{h_{a\parallel}B}=V_{2a}B+L_{ha}\fsa{B^2}$, де $V_{1,2a}=F\omega_{u,ha}/B$. Тоді усереднені паралельні баланси імпульсу і теплового потоку дають~\cite[с.~7, рівн.~73--78]{Seto2014}
\begin{equation}\label{eq:06-seto78}
  \fsa{3(\nabla_\parallel B)^2}
  \begin{pmatrix}\mu_{1a}&\mu_{2a}\\ \mu_{2a}&\mu_{3a}\end{pmatrix}
  \begin{pmatrix}L_{ua}\\ L_{ha}\end{pmatrix}
  =\sum_b\begin{pmatrix}l^{ab}_{11}&-l^{ab}_{12}\\ -l^{ab}_{21}&l^{ab}_{22}\end{pmatrix}
  \begin{pmatrix}V_{1b}B+L_{ub}\fsa{B^2}\\ V_{2b}B+L_{hb}\fsa{B^2}\end{pmatrix}
  +\begin{pmatrix}e_an_a\fsa{\Epar B}\\ 0\end{pmatrix}.
\end{equation}
Це стандартні неокласичні рівняння полоїдального обертання Хіршмана--Сігмара. Двовимірна модель їх відтворює після усереднення~\cite{Seto2014}, а \eqref{eq:06-utheta} є їх однокомпонентною границею. На LHD абсолютні полоїдальні швидкості і в іонному, і в електронному корені відповідають неокласичному рівню~\cite{Ida2001}. Стрибок $u_\theta$ при L--H переході в токамаках неокласикою не описується (п.~\ref{sec:06-shear}).

\section{Рівняння тороїдального імпульсу}\label{sec:06-tor}

Тороїдальну швидкість неокласика не фіксує, тому потрібне рівняння переносу. Singh та ін.~\cite{Singh2004} виводять його для краю в режимі Пфірша--Шлютера. Джерело нейтралів полоїдально асиметричне:
\begin{equation}\label{eq:06-neutrals}
  N_n=N_{0n}\bigl(1+\Delta_1\cos\theta\bigr),
\end{equation}
$\Delta_1>0$~--- напуск газу із зовнішнього боку, $\Delta_1<0$~--- із внутрішнього. Перезарядка вносить у \eqref{eq:06-mom} втрату імпульсу $M_{\mathrm{cx}}(\theta)$ і енергії $E_i(\theta)$, а іонізацію відкинуто. Далі розклад за полоїдальними гармоніками, усереднення по поверхні і границя $L_\perp\ll r\ll R$ з коловим перерізом дають одновимірне рівняння для густини кутового моменту~\cite[рівн.~45]{Singh2004}:
\begin{multline}\label{eq:06-singh}
  m_i\pa_t(n_iu_{\varphi i})+\nu_{\mathrm{cx}}m_in_iu_{\varphi i}+\pa_r\bigl[m_in_i\bigl(\langle\tilde u_E\tilde u_{\varphi}\rangle+u_{ri}u_{\varphi i}\bigr)\bigr]
  =\pa_r\biggl\{\eta_{2i}\biggl[\Bigl(1+1{,}67\,q\Delta_1\frac{\nu_{\mathrm{cx}}}{\nu_i}\frac{\Btor}{\Bpol}\Bigr)\frac{\pa\upar}{\pa r}\\
  +0{,}2\,q^2\frac{T_i}{e\Bpol}\Bigl(\frac{\pa\ln T_i}{\pa r}\Bigr)^{2}
  \Bigl(1-\frac{5{,}1\Delta_1}{q}\frac{\nu_{\mathrm{cx}}}{\nu_i}\frac{\Btor}{\Bpol}\Bigl(\rho_i\frac{\pa\ln T_i}{\pa r}\Bigr)^{-2}\Bigr)\biggr]\biggr\}.
\end{multline}
Тут $\eta_{2i}=4\eta_{1i}$, $\eta_{1i}=3p_i\nu_i/(10\Omc{i}^2)$~--- класична поперечна в'язкість, $\rho_i=c_i/\Omc{i}$, $c_i=\sqrt{T_i/m_i}$. Ліворуч стоять перезарядне гальмування, турбулентне напруження Рейнольдса $\langle\tilde u_E\tilde u_\varphi\rangle$ і неокласична конвекція. Праворуч~--- дифузія $\upar$ і неокласичне джерело від $\pa_r\ln T_i$ з поправкою асиметрії. При $\Delta_1=0$ відтворюється результат Claassen та ін.: $\pa_r\upar\approx-0{,}2q^2\frac{T_i}{e\Bpol}(\pa_r\ln T_i)^2$~\cite[рівн.~3, 47]{Singh2004}. Внесок іонної в'язкості відноситься до внеску нейтральної як $(\utor)_{\pi i}/(\utor)_{\pi n}\approx q^2/20\varepsilon\gtrsim1$, тож іонною в'язкістю нехтувати не можна~\cite[с.~2]{Singh2004}.

\begin{derivation}[оцінка швидкості біля LCFS~\cite{Singh2004}]
Нехтуємо ЛЧ \eqref{eq:06-singh} і вважаємо $\upar\approx\utor$ (бо $u_\theta\Bpol\ll\utor\Btor$). Тоді стаціонарне рівняння зводиться до рівності нулю виразу у квадратних дужках (рівн.~46). Покладемо масштаби градієнтів однаковими, $L_{T}=L_{T,U}$ (рівн.~48), і позначимо $(\rho_i)_p=(\Btor/\Bpol)\rho_i$. Отримаємо (рівн.~49)
$\bigl(1+1{,}67q\Delta_1\frac{\nu_{\mathrm{cx}}}{\nu_i}\frac{\Btor}{\Bpol}\bigr)\frac{\upar}{c_i}=-0{,}2q^2\frac{(\rho_i)_p}{L_T}\bigl[1-\frac{5{,}1\Delta_1}{q}\frac{\nu_{\mathrm{cx}}}{\nu_i}\frac{\Btor}{\Bpol}\frac{L_T^2}{(\rho_i)_p^2}\bigr]$.
Далі беремо типові для обох установок значення $q\frac{\nu_{\mathrm{cx}}}{\nu_i}\frac{\Btor}{\Bpol}\sim1$ і $q^2\frac{(\rho_i)_p}{L_T}\sim1$ та вважаємо $|\Delta_1|\ll1$. Тоді квадратна дужка дорівнює $1-5{,}1q^2\Delta_1$, і
\begin{equation}\label{eq:06-Uphi}
  |u_{\varphi i}|\approx0{,}2\,\bigl(1-5{,}1\,q^2\Delta_1\bigr)\,c_i .
\end{equation}
Із (46) знак \emph{ширу} $\pa_r\upar$ протилежний до знака квадратної дужки; сам знак $\upar$ локальний аналіз не визначає (у (48) неявно $\upar>0$, тож знак «$-$» у (49) джерела з (48) несумісний). Тому (51) подано за модулем.
\textbf{Увага:} у (45)--(46) джерела стоїть $\rho_i$, а в (49)~--- $(\rho_i)_p$; результат $1-5{,}1q^2\Delta_1$ отримано саме з $(\rho_i)_p$, з $\rho_i$ поправка була б у $(\Btor/\Bpol)^2\approx144$ рази більшою (TEXTOR, $\Btor/\Bpol\approx12$). Це ще одна внутрішня неузгодженість~\cite{Singh2004}.
\end{derivation}

\begin{corpusexample}[COMPASS-D і TEXTOR~\cite{Singh2004}]
Край COMPASS-D: $T_i=\SI{50}{еВ}$, $n_i=\SI{2e19}{м^{-3}}$, $\Bzero=\SI{1,2}{Тл}$, $\Ip=\SI{0,19}{МА}$, $R_0=\SI{0,56}{м}$, $a=\SI{0,17}{м}$, $c_i\approx\SI{7e4}{м/с}$, $\nu_i\approx\SI{4e4}{с^{-1}}$, $\Btor/\Bpol\approx5{,}4$ (так у джерелі), $\nu_{\mathrm{cx}}/\nu_i\sim10^{-2}$, $q_a\approx4{,}5$. TEXTOR (RI-мода): $T_i\approx\SI{50}{еВ}$, $\Bzero=\SI{2,25}{Тл}$, $\Ip=\SI{400}{кА}$, $R_0=\SI{1,75}{м}$, $a=\SI{0,46}{м}$, $q_a=3{,}4$, $c_i\approx\SI{5e4}{м/с}$, $\rho_i\approx\SI{4,46e-4}{м}$, $\Btor/\Bpol\approx12$, $\nu_{\mathrm{cx}}/\nu_i\approx2{,}3\cdot10^{-2}$. Без асиметрії \eqref{eq:06-Uphi} дає $\utor\sim0{,}2c_i\sim\SI{10}{км/с}$. Напуск із внутрішнього боку з $|\Delta_1|\sim2/(5q^2)$ дає $\sim\SI{30}{км/с}$, і при переході від зовнішнього напуску до внутрішнього шир $\pa_r\utor$ може змінити знак. \textbf{Застереження.} З балансу частинок ($\tau_{\mathrm{part}}\sim\SI{30}{мс}$, напуск $(0{,}1\text{--}1)\cdot10^{20}$~част./с) ті самі автори оцінюють $\Delta_1\sim(0{,}5\text{--}5)\cdot10^{-3}$. Для TEXTOR це в 7--70 разів менше за $2/(5q^2)\approx0{,}035$, тож «\SI{30}{км/с}» слід сприймати як якісний висновок (задача~3).
\end{corpusexample}

\begin{outofcorpus}[момент сил від NBI]
Пучок потужністю $P$ з енергією частинок $E_b$ і швидкістю $v_b$ вносить $P/E_b$ частинок за секунду з імпульсом $m_bv_b$ кожна. Нехай $R_{\tan}$~--- радіус дотику. Повний момент сил $T_{\mathrm{NBI}}=\frac{P}{E_b}m_bv_bR_{\tan}=\frac{2PR_{\tan}}{v_b}$: для D з енергією \SI{80}{кеВ} ($v_b\approx\SI{2,8e6}{м/с}$) і $R_{\tan}=\SI{2,5}{м}$ це $\approx\SI{1,8}{Н.м}$ на мегават. Радіальний профіль джерела $S_\varphi(r)$ повторює профіль осадження пучка (розд.~\ref{ch:heating}). У~\cite{Sigmar1987} це джерело $\vect{M}_j$, а гальмування описано як $-m_jn_j\nu_d\uvec_j$ з експериментальною $\nu_d$.
\end{outofcorpus}

\section{Власне обертання і недифузійний потік імпульсу}\label{sec:06-intrinsic}

Замикання для напруження Рейнольдса в \eqref{eq:06-singh} Ida~\cite{Ida2001} отримує з досліду JFT-2M. Потік тороїдального імпульсу складається з дифузійного і недифузійного (off-diagonal) доданків. Формули в статті немає, тож нижче наш запис, а коефіцієнт $\alpha$ з огляду перейменовано на $C_T$, бо α зайнята α-частинками:
\begin{equation}\label{eq:06-Pi}
  \Pitor=-\chi_{\varphi}m_in_i\frac{\pa\utor}{\pa r}+\Pi_{\mathrm{res}},\qquad \Pi_{\mathrm{res}}\propto\frac{\pa T_i}{\pa r}.
\end{equation}
Відношення недифузійного коефіцієнта до дифузійного~--- 0,1--0,3~\cite[с.~3--5]{Ida2001}. На рис.~3(b) статті воно спадає від $\approx0{,}15$ при \SI{155}{кА} до $\approx0{,}05$ при \SI{305}{кА}. Нормування осі в короткій статті не пояснено, тому числа порівнюйте лише за порядком. У практикумі використано наше безрозмірне нормування $\Pi_{\mathrm{res}}=-C_T\chi_{\varphi}m_in_ic_i\,\pa_r\ln T_i$.

\begin{corpusexample}[JFT-2M, CHS, LHD~\cite{Ida2001}]
\emph{JFT-2M}~\cite[с.~3--5]{Ida2001} ($\Ip=\SI{243}{кА}$). Схема досліду: при $t=\SI{750}{мс}$ пучки перемикають між co- і counter-інжекцією, зберігаючи однаковий поглинутий імпульс. Результати: (1)~у центрі ($R=\SI{1,35}{м}$) counter-NBI дає приблизно вдвічі швидше обертання, ніж co-NBI; (2)~у шарі $R=1{,}4$--\SI{1,5}{м} градієнт швидкості відрізняється в 2--3 рази; (3)~ця різниця зростає лінійно з $\pa_rT_i$ (при $\rho=0{,}46$); (4)~недифузійний коефіцієнт зменшується зі зростанням $\Ip$ (можливо, як $1/\Ip$ чи $1/\Bpol$). Втратами на орбітах це не пояснити: ефект є і на великих токамаках. \emph{CHS}~\cite[с.~4--6]{Ida2001}: у режимі гарячих електронів (друга ЕЦ-гармоніка \SI{0,11}{МВт}, NBI \SI{0,7}{МВт}, $n_e<\SI{0,7e19}{м^{-3}}$, $T_e$ зростає з 0,2 до \SI{2}{кеВ}) центральний спонтанний потік до \SI{5e4}{м/с} перевершує і розвертає обертання від NBI; при $\rho=0{,}3$ маємо $\Er=\SIrange{10}{20}{кВ/м}$, полоїдальний потік \SIrange{10}{20}{км/с}. \emph{LHD}~\cite[с.~2--3]{Ida2001}: за $n_e=(0{,}4\text{--}1{,}0)\cdot10^{19}$~м$^{-3}$ сам NBI переводить край ($\rho>0{,}85$) з іонного кореня в електронний. В електронному корені $\Er$ швидко росте зі зменшенням густини до $+\SI{15}{кВ/м}$; в іонному ($n_e=(1{,}0\text{--}3{,}0)\cdot10^{19}$~м$^{-3}$) воно плавно доходить до $-\SI{5}{кВ/м}$.
\end{corpusexample}

\textbf{Напрям.} Пояснення Іди спирається на симетрію поля~\cite[с.~6--7]{Ida2001}. Спонтанний потік обирає той напрям на поверхні, у якому $B$ змінюється найменше: там паралельна в'язкість, що гальмує потік, найслабша. У токамаку таким напрямом завдяки осесиметрії є тороїдальний, тому потік, керований $\grad T_i$, тороїдальний і паралельний дрейфу $\fsa{\Er\times\Bpol}$ (протиструмовий, підсилює від'ємне $\Er$). У геліотроні найменша зміна $B$ відбувається вздовж гвинтової лінії, тож потік гвинтовий, і його тороїдальна проєкція може бути антипаралельною. Рівняння лінійне, тож обертання від пучка і власне обертання додаються: $u_{\mathrm{co,ctr}}=\pm u_{\mathrm{NBI}}+u_{\mathrm{int}}$. Відношення $|u_{\mathrm{ctr}}|/|u_{\mathrm{co}}|=2$ означає $|u_{\mathrm{int}}|=u_{\mathrm{NBI}}/3$ (наш висновок; перевірка в практикумі).

\section{Шир \texorpdfstring{$\ExB$}{E×B} і турбулентність}\label{sec:06-shear}

Критерій придушення турбулентності Singh та ін. записують як~\cite[рівн.~1]{Singh2004}
\begin{equation}\label{eq:06-shear}
  \gamma_E=\Bigl|\Delta_r\frac{\pa\omega_{\mathrm D}}{\pa r}\Bigr|\ge\gamma_T,\qquad \omega_{\mathrm D}=-\frac{k_\theta\Er}{B},
\end{equation}
де $\omega_{\mathrm D}$~--- доплерівський зсув частоти від руху $\ExB$, $k_\theta$~--- полоїдальне хвильове число, $\Delta_r$~--- радіальна ширина лінійної моди, $\gamma_T$~--- інкремент фонової нестійкості (ITG, резистивні балонні моди). Разом з \eqref{eq:06-Erneo} це дає два «важелі» для керування $\gamma_E$: $\pa_rp_i$ і $\pa_r\utor$. Саме так асиметричний напуск змінює поріг L--H~\cite{Singh2004}. Експериментально на транспортних бар'єрах темп зсуву $\ExB$ справді перевищує $\gamma_T$~\cite{Ida2001}.

\begin{outofcorpus}[шир у тороїдальній геометрії і множник для $\chi_s$]
Інваріантна форма темпу зсуву: $\gamma_{E\times B}=\frac{(R\Bpol)^2}{B}\frac{\pa}{\pa\psi}\Bigl(\frac{\Er}{R\Bpol}\Bigr)$ (Хам--Беррелл), у \eqref{eq:06-shear} це $\gamma_E$. Шир твердотільного обертання $\Er\propto R\Bpol$ турбулентність не пригнічує. У 1.5D-кодах вплив ширу вводять множником $\chi_s\to\chi_s/[1+(\gamma_E/\gamma_T)^2]$ або обрізанням $\gamma_T\to\max(\gamma_T-\gamma_E,0)$; $\gamma_T$ беруть з гірокінетики чи квазілінійної моделі.
\end{outofcorpus}

\textbf{H-мода.} При переході на краю стрибком з'являється полоїдальний потік в електронному діамагнітному напрямку з максимумом біля сепаратриси. Поза сепаратрисою плазма обертається в іонному напрямку, тож біля неї $\Er$ має «яму» з великим широм~\cite{Ida2001}. Критичної зіткненнєвості $\nu_{*i}$ для переходу не виявлено, а його величина не залежить від полоїдального ларморового радіуса~\cite{Ida2001}. У геліотронах, навпаки, полоїдальний потік перемикається з електронного напрямку в іонний, бо велике $\Er$ зменшує гофрувальні втрати~\cite{Ida2001}.

\textbf{Зональні потоки.} Глобальне гірокінетичне моделювання GTC у геометрії DIII-D~\cite[с.~15--18]{Wang2006} дає таку картину. Нестійкість ITG збуджується в зоні $0{,}42<r<0{,}76$ і породжує радіально витягнуті «стримери». Самозбуджені зональні потоки $\ExB$ розривають їх на дрібніші радіальні масштаби, після чого турбулентність розтікається в лінійно стійкі області з порівнянною інтенсивністю (нелокальний перенос). У беззіткненнєвій плазмі турбулентність і зональний потік регулюють одне одного і дають сплески, повільніші за коливання GAM; без зональних потоків сплесків немає. Штучні зональні потоки можуть збуджувати турбулентність (з порогом і запізненням), але її енергія на два порядки менша за енергію зональних потоків, тож як механізм насичення це надто слабко~\cite{Wang2006}.

\textbf{Домішки.} Вплив швидкого обертання на неокласичну дифузію важких домішок розглянуто у п.~\ref{sec:05-imp}, формула~\eqref{eq:05-wong}. На PLT при counter-NBI у ядро потрапляло в 2--3 рази більше домішок, ніж при co-NBI~\cite[с.~3]{Wong1987}. Отже, напрям обертання є параметром керування домішками.

\section{Потоки в шарі зіскрібання}\label{sec:06-sol}

Поза сепаратрисою силові лінії впираються в мішені, і на вході в шар діє критерій Бома $\mathcal{M}=u_\parallel/c_s\ge1$, $c_s=\sqrt{(ZT_e+\gamma_{\mathrm a}T_i)/m_i}$ ($\gamma_{\mathrm a}$~--- показник адіабати іонів)~\cite[с.~24]{AhoMantila2011}. Потенціал плазми відносно мішені дорівнює падінню в шарі й передшарі, $\Phi\approx3{,}5\,T_e/e$ (розд.~\ref{ch:wall}, \eqref{eq:08-Vsh}). Звідси наша оцінка $\Er^{\mathrm{SOL}}\approx-3{,}5\,\pa_rT_e/e>0$: у SOL поле додатне, а під сепаратрисою від'ємне, тож на сепаратрисі воно змінює знак. Корпус описує цей механізм якісно: радіальний градієнт $T_e$ створює $\Er$, дрейф $\Er\times\Bvec$ іде полоїдально, а асиметрію тиску між мішенями компенсує зворотний паралельний потік~\cite{AhoMantila2011}. Паралельний потік Пфірша--Шлютера компенсує дивергенцію вертикального $\grad B$-дрейфу.

\begin{corpusexample}[{ASDEX Upgrade, SOLPS5.0~\cite[с.~61--64]{AhoMantila2011}}]
\emph{Прямий дрейф} ($\Bvec\times\grad B$ до дивертора): у головному SOL виміряно паралельні потоки з $\mathcal{M}\sim0{,}5$, а двовимірні коди дають у 2--10 разів менше. Можлива причина~--- заниження розрахункового $\Er$ у SOL; якщо підставити виміряне $\Er$, потоки Пфірша--Шлютера зростають. Для L-моди низької густини SOLPS5.0 з дрейфами відтворює профілі на LFS, проте паралельні іонні потоки все одно в 2--3 рази нижчі за виміряні. \emph{Обернене поле}: згода кількісна. \emph{Домішки}: перенесення вуглецю в зовнішньому диверторі задають дрейф $\ExB$ і тертя об потік плазми; при оберненні поля «хвости» осадження змінюють напрям, а виміряна ефективність локального осадження 24--32\,\%~\cite[с.~63]{AhoMantila2011} падає приблизно вдвічі.
\end{corpusexample}

\section{Вимірювання обертання і \texorpdfstring{$\Er$}{Er}}\label{sec:06-meas}

\begin{outofcorpus}[CXRS і HIBP]
\emph{Спектроскопія перезарядки (CXRS)}: нейтрали пучка віддають електрон повністю іонізованій домішці (C$^{6+}$, Ne$^{10+}$), і збуджений іон випромінює. Доплерівський зсув лінії дає $u_\varphi$ чи $u_\theta$ домішки, ширина~--- $T_I$, а перетин з пучком~--- локальність. $\Er$ отримують з \eqref{eq:06-Er} для домішки, де доданок тиску малий, як $1/Z_I$. \emph{Зонд пучком важких іонів (HIBP)} вимірює локальний потенціал за енергією вторинних іонів, тобто $\Er$ без балансу сил.
\end{outofcorpus}

У CHS $\Er$, відновлене зі спектроскопічно виміряних потоків, збігається з виміром HIBP~\cite{Ida2001}. Це пряма перевірка \eqref{eq:06-Er}. Для густої плазми Марценюк~\cite[с.~119--126]{Martsenyuk2025} розвиває \emph{мікрохвильову рефракцію}. Звичайна хвиля частоти $\omega_0$ відбивається від шару з критичною густиною $N_{\mathrm{cr}}=\varepsilon_0m_e\omega_0^2/e^2$, а неоднорідності цього шару модулюють сигнал на приймальних рупорах. Зі взаємної кореляції сигналів двох антен, рознесених на кут $\Delta\varphi_{\mathrm a}$, кутову частоту знаходять за часовим зсувом максимумів: $\omega=\Delta\varphi_{\mathrm a}/\Delta\tau$. На установці MAKET (комірка Пеннінга, лінійна, \emph{не} тороїдальна) схема з чотирма антенами виявила чотири жолобки густини з радіальним розміром $\approx\SI{1}{см}$ і азимутальними зсувами $47^\circ$, $53^\circ$, $53^\circ$, $207^\circ$. Плазма обертається за годинниковою стрілкою з $\omega\approx\SI{1,16e4}{рад/с}$ у схрещених полях $\ExB$~\cite{Martsenyuk2025}. Метод цікавий як незалежна від спектроскопії діагностика, але в корпусі його не застосовано до токамака.

\section{Синтез: ієрархічна модель потоку для проєкту}\label{sec:06-model}

Нижче~--- замкнена 1.5D-система (табл.~\ref{tab:06-model}; циліндричне наближення $r$, $\Bpol=r\Btor/(qR_0)$; у формованій геометрії $\frac1r\pa_r r\to\frac{1}{V'}\pa_rV'\fsa{|\grad r|^2}$, розд.~\ref{ch:geometry}). Невідомі: $n_e,n_i,T_e,T_i,\psi$ (розд.~\ref{ch:transport}), $\utor$, $u_\theta$, $\Er$.
\begin{enumerate}[label=(\roman*)]
\item \emph{Геометрія}: $\GSop\psi=-\mu_0R^2p'-FF'$ дає $q(r)$, $\Bpol$, $V'$ (розд.~\ref{ch:equilibrium}).
\item \emph{1.5D-транспорт} $n,T_{e,i},\psi$ з $\chi_s$, помноженими на $f_E$ з кроку (vi).
\item \emph{Тороїдальний імпульс}:
\begin{equation}\label{eq:06-model-mom}
  \begin{gathered}
  m_i\pa_t(n_i\utor)+\frac1r\pa_r\bigl(r\Pitor\bigr)=-\nu_{\mathrm{cx}}m_in_i\utor+S_{\mathrm{NBI}}+S_{\mathrm{neo}},\\
  \Pitor=-\chi_{\varphi}m_in_i\pa_r\utor-C_T\chi_{\varphi}m_in_ic_i\pa_r\ln T_i .
  \end{gathered}
\end{equation}
\item \emph{Полоїдальний потік}: $u_\theta=\mathcal{K}(\nustar)\,\pa_rT_i/(e\Btor)$ \eqref{eq:06-utheta} або \eqref{eq:06-seto78} для кількох сортів.
\item \emph{Радіальне поле}: \eqref{eq:06-Er} або \eqref{eq:06-Erneo}.
\item \emph{Зворотний зв'язок}: $\gamma_E=|\pa_r(\Er/B)|$ (при $k_\theta\Delta_r\sim1$), $f_E=[1+(\gamma_E/\gamma_T)^2]^{-1}$.
\item \emph{Граничні умови}: $\pa_r\utor(0)=0$; на LCFS $\utor(a)=u_a$ (оцінка \eqref{eq:06-Uphi}) або SOL-модель з критерієм Бома.
\end{enumerate}

\begin{table}[ht]
\centering\small
\caption{Звідки береться кожен член моделі п.~\ref{sec:06-model}. К~--- корпус, П~--- поза корпусом, Н~--- наша нормалізація/оцінка; ПШ~--- режим Пфірша--Шлютера}\label{tab:06-model}
\begin{tabularx}{\textwidth}{@{}p{4.1cm}Xc@{}}
\toprule
Член & Джерело і межі застосовності & Статус\\
\midrule
Баланс сил \eqref{eq:06-Er}, \eqref{eq:06-seto} & \cite{Seto2014,Singh2004,Ida2001}; нульовий порядок за δ & К\\
Форма потоку \eqref{eq:06-uform} & \cite{Seto2014} (рівн.~53), \cite{Sigmar1987} при $u\sim v_{\mathrm{th}}$ & К\\
Рівняння для $\omega$ \eqref{eq:06-omega}, біфуркація & \cite{Sigmar1987}; лише ПШ, ізотермічні поверхні; стійкість коренів & К / Н\\
$u_\theta$: ПШ ($-1{,}83$), плато, банан & \cite{Singh2004} (числа); виведення бананового & К / П\\
Багатосортова неокласика \eqref{eq:06-seto78} & \cite{Seto2014}; $\mu_{ai}$, $l^{ab}_{ij}$~--- Хіршмана--Сігмара & К / П\\
$\nu_{\mathrm{cx}}$, $S_{\mathrm{neo}}(\Delta_1)$ \eqref{eq:06-singh} & \cite{Singh2004}; край, ПШ, коловий переріз & К\\
$S_{\mathrm{NBI}}$ & стандартна формула моменту сил & П\\
$\chi_{\varphi}$, $\Pi_{\mathrm{res}}\propto\pa_rT_i$ & \cite{Ida2001}: відношення 0,1--0,3, $\propto1/\Ip$; нормування $C_T$ & К / Н\\
Критерій $\gamma_E\ge\gamma_T$ & \cite{Singh2004,Ida2001}; форма $f_E$, $\gamma_T$ & К / П\\
Зональні потоки, нелокальність & \cite{Wang2006}~--- якісно (DIII-D, ITG) & К\\
Межа: Бом, SOL-потоки & \cite{AhoMantila2011}; $\Er^{\mathrm{SOL}}$ з падіння в шарі & К / Н\\
Валідація & \cite{Ida2001,Singh2004,Wong1987}: порядки $\Er$, $u_\theta$, $\utor$ & К\\
\bottomrule
\end{tabularx}
\end{table}

\textbf{Цілі валідації} (порядки величин): $|\Er|\sim5$--20~кВ/м (LHD, CHS), $u_\theta\sim10$--20~км/с (CHS), $\utor$ до \SI{5e4}{м/с} (CHS), $\omtor(0)\sim\SI{2e5}{рад/с}$ (модельні профілі TFTR~\cite{Wong1987}), $\utor\sim10$--30~км/с біля LCFS (оцінка~\cite{Singh2004}), $\mathcal{M}\sim0{,}5$ у SOL~\cite{AhoMantila2011}. \textbf{Прогалина JET.} У корпусі немає жодного виміру обертання, $\Er$ чи потоків у SOL на JET. JET фігурує лише як сценарій пуску~\cite{Hoppe2022}, рядок параметрів для W~\cite{vanVugt2019}, параметри розряду~\cite{Fukuyama1992} і профілі $T_{e,i}$ DTE1~\cite{Tyshchenko2021}. Тому модель для JET можна перевірити лише за порядками величин з інших установ. Для кількісної валідації потрібні дані CXRS JET і розрахунки JINTRAC/JETTO поза корпусом. Усі теоретичні ланки \eqref{eq:06-singh}, \eqref{eq:06-omega} виведено для режиму Пфірша--Шлютера, великого аспектного відношення і колового перерізу. Для ядра JET (банановий режим, $\kappa\approx1{,}7$) їх треба замінити стандартною неокласикою.

\begin{practicum}[1D тороїдальний імпульс і $\Er$: новий модуль \texttt{flows\_1d.py}]
Модуля потоків у коді проєкту немає. Створіть \texttt{Our~try/\allowbreak 03-deep-dives/\allowbreak D4-flows/\allowbreak flows\_1d.py} (numpy/scipy).
\textbf{Вхідні дані} (TFTR, модельні профілі~\cite{Wong1987}): $a=\SI{0,82}{м}$, $\Btor=\SI{4}{Тл}$, D-плазма; $n_i=\SI{2e19}{м^{-3}}/(1+4x^2)$, $T_i=\SI{10}{кеВ}/(1+9x^2)$, $q=1{,}1(1+6{,}5455x^2)$, $\omtor=\SI{2e5}{рад/с}/(1+4x^2)$, $\utor=\omtor R_0$, $x=r/a$. Припущення поза корпусом: $R_0=\SI{2,5}{м}$, $\chi_{\varphi}=\SI{1}{м^2/с}$.
\textbf{A. Баланс сил.} При $r=a/2$ обчисліть три доданки \eqref{eq:06-Er} з $u_\theta$ за \eqref{eq:06-utheta}. Очікувано: $q=2{,}90$, $\Bpol=\SI{0,226}{Тл}$, тиск $\approx-\SI{17,9}{кВ/м}$, $\utor\Bpol\approx+\SI{56,6}{кВ/м}$, полоїдальний $\approx+\SI{12,2}{кВ/м}$ ($\mathcal{K}=1{,}17$), разом $\Er\approx\SI{51}{кВ/м}$.
\textbf{B. Стаціонарне \eqref{eq:06-model-mom}} (скінченні об'єми, $N=400$, $\pa_r\utor(0)=0$, $\utor(a)=0$). \emph{Перевірки}: при сталих $n$, $\chi_{\varphi}$, $\nu_{\mathrm{cx}}=\SI{10}{с^{-1}}$ і сталому $S$ розв'язок має збігтися з $\utor=\frac{S}{m_in\nu_{\mathrm{cx}}}\bigl[1-\frac{I_0(r/\lambda)}{I_0(a/\lambda)}\bigr]$, $\lambda=\sqrt{\chi_{\varphi}/\nu_{\mathrm{cx}}}$ (похибка $\sim10^{-5}$), а при $\nu_{\mathrm{cx}}=0$~--- з параболою $S(a^2-r^2)/(4\chi_{\varphi}m_in)$. \emph{Робочий випадок}: $\nu_{\mathrm{cx}}=\SI{2e3}{с^{-1}}\,\ee^{(r-a)/(0{,}05\,\text{м})}$ (порядок $\nu_{\mathrm{cx}}/\nu_i\sim10^{-2}$ на краю~\cite{Singh2004}), $S=S_0\ee^{-(r/0{,}4\,\text{м})^2}$. Підберіть $S_0$ так, щоб $\utor(0)=\omtor(0)R_0=\SI{5e5}{м/с}$. Очікувано: $S_0\approx\SI{0,235}{Н/м^3}$, момент сил $\approx\SI{4,6}{Н.м}$ (порівняйте з NBI, $\approx\SI{1,8}{Н.м/МВт}$), $\tau_\varphi=L_\varphi/T\approx\SI{0,13}{с}$, $\omtor(a/2)\approx\SI{1,3e5}{рад/с}$ (у Wong--Cheng $\SI{1e5}{рад/с}$). Без перезарядки $\utor(0)\approx\SI{7,1e5}{м/с}$.
\textbf{C. Власне обертання.} Увімкніть $C_T$ і порівняйте $S_0\to\pm S_0$. При $C_T=0{,}1$ і $0{,}3$ відношення $|u_{\mathrm{ctr}}(0)|/|u_{\mathrm{co}}(0)|\approx1{,}4$ і $3{,}1$; значення 2 (JFT-2M) дає $C_T\approx0{,}19$ (збіг лише за порядком: установки й нормування різні).
\textbf{D.} Підставте розв'язок у \eqref{eq:06-Er}, обчисліть $\gamma_E(r)$ і множник $f_E$ для $\gamma_T=\SI{1e5}{с^{-1}}$ (задане). Факультативно візьміть $q(\psiN)$ DIII-D з логу \texttt{run\_bake.py} (\texttt{tokamak-3d-viz}, розд.~\ref{ch:particles})~\cite{ProjViz}.
\end{practicum}

\subsection*{Контрольні запитання}
\begin{enumerate}
  \item Виведіть \eqref{eq:06-Er} з \eqref{eq:06-seto}. Як зміняться знаки, якщо θ відлічувати проти годинникової стрілки? Чому для відновлення $\Er$ вигідніше брати домішку, а не основні іони?
  \item Чому потік на поверхні задають саме дві функції потоку \eqref{eq:06-uform}? Покажіть, що $u_\theta\propto\Bpol/n$ і що при $\omega_\psi=0$ потік суто паралельний $\Bvec$.
  \item Чому неокласична в'язкість фіксує полоїдальне обертання, але не тороїдальне? Як це пов'язано з осесиметрією~\cite{Seto2014,Ida2001} і з розділенням $\tau_p\ll\tau_T$~\cite{Sigmar1987}?
  \item Що змінює знак $\mathcal{K}$ між банановим режимом і режимом Пфірша--Шлютера і як це впливає на внесок $\pa_rT_i$ у $\Er$ \eqref{eq:06-Erneo}?
  \item Які спостереження JFT-2M свідчать про недифузійний потік імпульсу? Чому напрям власного обертання в токамаку і геліотроні протилежний~\cite{Ida2001}?
  \item Чому шир твердотільного обертання не пригнічує турбулентність? Чому, за~\cite{Wang2006}, перенесення енергії від зональних потоків до турбулентності не насичує зональні потоки?
\end{enumerate}

\subsection*{Задачі}
\begin{enumerate}
  \item (Ida~\cite{Ida2001}; Марценюк~\cite{Martsenyuk2025}.) а) У CHS $\Er\approx u_\theta\Btor$. Яке поле $\Btor$ випливає з $\Er=10$--20~кВ/м і $u_\theta=10$--20~км/с? б) На MAKET $\omega=\SI{1,16e4}{рад/с}$. Знайдіть період обертання і час $\Delta\tau$ між проходами одного жолобка повз антени, рознесені на $90^\circ$, а також між жолобками зі зсувом $47^\circ$ повз одну антену.
  \item (Singh та ін.~\cite{Singh2004}; край TEXTOR.) $T_i=\SI{50}{еВ}$, $\Btor=\SI{2,25}{Тл}$, $\Btor/\Bpol=12$, $q=3{,}4$, $\rho_i=\SI{4,46e-4}{м}$. Нехай $q^2(\rho_i)_p/L_T=1$ і $L_n=L_T$. Знайдіть $L_T$, $u_\theta$ у режимі ПШ і модулі трьох доданків \eqref{eq:06-Er} при $|\utor|=\SI{10}{км/с}$.
  \item (Там само.) За \eqref{eq:06-Uphi} з $c_i=\SI{5e4}{м/с}$ і $q=3{,}4$ обчисліть $|\utor|$ для внутрішнього напуску при $|\Delta_1|=2/(5q^2)$, $5\cdot10^{-3}$, $5\cdot10^{-4}$ і для зовнішнього при $\Delta_1=5\cdot10^{-3}$. Наскільки висновок про \SI{30}{км/с} чутливий до $\Delta_1$?
  \item (Sigmar та ін.~\cite{Sigmar1987}.) Для \eqref{eq:06-omega} з $C_1>0$ знайдіть стаціонарні корені, покажіть, що стійкий лише $\omega_+$, і знайдіть час релаксації. Що відбувається при $C_2^2=4C_1C_3$? Перепишіть $D$ і зв'язок $\Er=R\Bpol\omega/c$ з СГС у СІ.
\end{enumerate}
